\documentclass[11pt]{article}
\usepackage{amsmath,amssymb,amsthm}
\newtheorem{lemma}{Lemma}
\newtheorem{proposition}{Proposition}
\newtheorem{remark}{Remark}
\begin{document}

\section*{Detailed derivation of the block conditions for \(X,Y\)}

\subsection*{Setup}
Let \(T,T'\in B(H)\) be isometries with \(TT'=T'T\) (commuting only).
Set
\[
D:=D_{T^*}=(I-TT^*)^{1/2}=P_{\ker T^*}.
\]
Halmos unitary extension of \(T\) on \(H\oplus H\):
\[
U_0=\begin{pmatrix} T & D \\ 0 & -T^* \end{pmatrix},
\qquad
U_0^*=\begin{pmatrix} T^* & 0 \\ D & -T \end{pmatrix}.
\]
Let \(\mathcal{M}=\bigvee_{n\in\mathbb{Z}}U_0^n(H\oplus 0)\) and \(U=U_0\big|_{\mathcal{M}}\).
Define \(\widetilde{T}'\) on \(\mathcal{M}\) by
\begin{equation}
\label{eq:def}
\widetilde{T}'(U^n h)=U^n(T'h)\qquad(n\in\mathbb{Z},\,h\in H).
\end{equation}
Then \(\widetilde{T}'|_H=T'\) and \(\widetilde{T}'U=U\widetilde{T}'\).

Identify \(h\in H\) with \(\begin{pmatrix}h\\ 0\end{pmatrix}\in\mathcal{M}\).

\bigskip
%% ============================================================
\section*{Part A --- Block shape \(\displaystyle\widetilde{T}'=\begin{pmatrix}T'&X\\ 0&Y\end{pmatrix}\)}

From~\eqref{eq:def} with \(n=0\),
\[
\widetilde{T}'\begin{pmatrix}h\\ 0\end{pmatrix}
=
\begin{pmatrix}T'h\\ 0\end{pmatrix}.
\]
So \(H\oplus 0\) is invariant under \(\widetilde{T}'\), and the action there is \(T'\).
Any bounded operator on (a reducing piece of) \(H\oplus H\) with this property has the block form
\begin{equation}
\label{eq:block}
\boxed{
\widetilde{T}'
=
\begin{pmatrix}
T' & X \\
0 & Y
\end{pmatrix}
}
\end{equation}
for unique \(X,Y\in B(H)\) (defined on the second-summand part of \(\mathcal{M}\)). Explicitly,
\begin{equation}
\label{eq:action}
\widetilde{T}'\begin{pmatrix}x\\ y\end{pmatrix}
=
\begin{pmatrix}
T'x+Xy\\
Yy
\end{pmatrix}.
\end{equation}
The bottom-left zero is exactly the statement that \(\widetilde{T}'\) extends \(T'\)
(first slot invariant, no feed from \(H\) into the second slot).

\bigskip
%% ============================================================
\section*{Part B --- Values on the defect: what \(X,Y\) do to \(\ker T^*\)}

\begin{lemma}
For every \(h\in H\),
\begin{equation}
\label{eq:Ustar}
U^*\begin{pmatrix}h\\ 0\end{pmatrix}
=
\begin{pmatrix}T^*h\\ Dh\end{pmatrix}.
\end{equation}
In particular, if \(d\in\ker T^*\), then \(T^*d=0\), \(Dd=d\), so
\begin{equation}
\label{eq:past}
\begin{pmatrix}0\\ d\end{pmatrix}
=
U^*\begin{pmatrix}d\\ 0\end{pmatrix}.
\end{equation}
\end{lemma}

\begin{proposition}[correct defect formulae]
For every \(d\in\ker T^*\),
\begin{equation}
\label{eq:XY-d}
\boxed{
\widetilde{T}'\begin{pmatrix}0\\ d\end{pmatrix}
=
\begin{pmatrix}
T^*T'd\\
D\,T'd
\end{pmatrix}.
}
\end{equation}
Comparing with~\eqref{eq:action},
\begin{equation}
\label{eq:XY-on-d}
\boxed{Xd=T^*T'd,\qquad Yd=D\,T'd.}
\end{equation}
\end{proposition}

\begin{proof}
Use~\eqref{eq:past} and commutation \(\widetilde{T}'U^*=U^*\widetilde{T}'\):
\begin{align*}
\widetilde{T}'\begin{pmatrix}0\\ d\end{pmatrix}
&=
\widetilde{T}'\,U^*\begin{pmatrix}d\\ 0\end{pmatrix}
=
U^*\,\widetilde{T}'\begin{pmatrix}d\\ 0\end{pmatrix}
=
U^*\begin{pmatrix}T'd\\ 0\end{pmatrix}
=
\begin{pmatrix}T^*(T'd)\\ D(T'd)\end{pmatrix}.
\end{align*}
On the other hand, by~\eqref{eq:action},
\[
\widetilde{T}'\begin{pmatrix}0\\ d\end{pmatrix}
=
\begin{pmatrix}Xd\\ Yd\end{pmatrix}.
\]
Equating components yields~\eqref{eq:XY-on-d}.
\end{proof}

\begin{remark}[important correction]
One does \emph{not} in general have \(Xd=0\). That would require \(T^*T'd=0\), i.e.\
\(T'd\in\ker T^*\), i.e.\
\[
\langle T'd,\,Th\rangle=\langle d,\,{T'}^*Th\rangle=0\quad(\forall h),
\]
which holds automatically for all \(d\in\ker T^*\) if and only if
\({T'}^*T=T{T'}^*\) on \(H\) --- equivalent to double commutativity
\(T^*T'=T'T^*\). In the merely commuting case, \(Xd=T^*T'd\) may be nonzero.
(The identity \(Yd=DT'd\) is always correct.)
\end{remark}

\begin{remark}[norm check]
For \(d\in\ker T^*\),
\[
\|Xd\|^2+\|Yd\|^2
=
\|T^*T'd\|^2+\|DT'd\|^2
=
\|T'd\|^2
=
\|d\|^2,
\]
so~\eqref{eq:XY-on-d} is compatible with isometry of \(\widetilde{T}'\).
\end{remark}

\bigskip
%% ============================================================
\section*{Part C --- Commutation \(\widetilde{T}'U=U\widetilde{T}'\)}

Compute, with~\eqref{eq:block},
\begin{align*}
\widetilde{T}'U_0
&=
\begin{pmatrix}T'&X\\ 0&Y\end{pmatrix}
\begin{pmatrix}T&D\\ 0&-T^*\end{pmatrix}
=
\begin{pmatrix}
T'T & T'D-XT^* \\
0 & -YT^*
\end{pmatrix},
\\[1mm]
U_0\widetilde{T}'
&=
\begin{pmatrix}T&D\\ 0&-T^*\end{pmatrix}
\begin{pmatrix}T'&X\\ 0&Y\end{pmatrix}
=
\begin{pmatrix}
TT' & TX+DY \\
0 & -T^*Y
\end{pmatrix}.
\end{align*}

\begin{proposition}
On \(\mathcal{M}\), \(\widetilde{T}'U=U\widetilde{T}'\) holds if and only if
\begin{equation}
\label{eq:comm}
\boxed{
YT^*=T^*Y,
\qquad
T'D-XT^*=TX+DY.
}
\end{equation}
\end{proposition}

\begin{proof}
The \((1,1)\)-blocks agree by \(TT'=T'T\). Equating \((2,2)\) gives
\(YT^*=T^*Y\). Equating \((1,2)\) gives \(T'D-XT^*=TX+DY\).
\end{proof}

\paragraph{Meaning of \(YT^*=T^*Y\).}
\(Y\) intertwines \(T^*\) with itself. So \(Y\) preserves the spectral / Wold structure
of \(T^*\) (it maps \(\ker T^*\) into itself up to the unitary part, etc.).

\paragraph{Meaning of \(T'D-XT^*=TX+DY\).}
The Halmos operator \(U_0\) mixes the two slots through the off-diagonal \(D\).
For \(\widetilde{T}'\) to commute with that mixing, the ``naive'' action \(T'D\) on the
defect must be corrected by the unknown \(X\). Rewrite the identity as
\begin{equation}
\label{eq:X-solve}
X T^* + T X
=
T'D-DY.
\end{equation}
On \(\operatorname{Ran}D=\ker T^*\) one has \(T^*|_{\ker T^*}=0\), so~\eqref{eq:X-solve}
reduces to a formula for \(TX\) in terms of known data once \(Y\) is known on the defect;
together with~\eqref{eq:XY-on-d} this determines the germ of \(X,Y\) on the wandering
subspace. The inductive definition~\eqref{eq:def} then propagates the values to all of
\(\mathcal{M}\) via powers of \(U\).

\bigskip
%% ============================================================
\section*{Part D --- Isometry \(\widetilde{T}'{}^*\widetilde{T}'=I\)}

\begin{proposition}
\(\widetilde{T}'\) in~\eqref{eq:block} is an isometry iff
\begin{equation}
\label{eq:isom}
\boxed{
{T'}^*X=0,
\qquad
X^*X+Y^*Y=I.
}
\end{equation}
\end{proposition}

\begin{proof}
\[
{\widetilde{T}'}^*\widetilde{T}'
=
\begin{pmatrix}{T'}^*&0\\ X^*&Y^*\end{pmatrix}
\begin{pmatrix}T'&X\\ 0&Y\end{pmatrix}
=
\begin{pmatrix}
I & {T'}^*X \\
X^*T' & X^*X+Y^*Y
\end{pmatrix}.
\]
\end{proof}

\paragraph{Meaning of \({T'}^*X=0\).}
\(\operatorname{Ran}X\subset\ker{T'}^*=H\ominus T'H\): the off-diagonal \(X\) lands in the
defect space of \(T'\).

\paragraph{Meaning of \(X^*X+Y^*Y=I\).}
On the second Halmos slot, the column \(\begin{pmatrix}X\\ Y\end{pmatrix}\) is an isometry:
it splits the second slot between the two output slots without losing norm
(cf.\ the norm check after~\eqref{eq:XY-on-d}).

\bigskip
%% ============================================================
\section*{Part E --- How the pieces fit together}

\begin{enumerate}
\item \textbf{Definition~\eqref{eq:def}} produces a concrete isometry \(\widetilde{T}'\) commuting
with \(U\) and extending \(T'\). Existence is free; no double commutativity is used
(only \(TT'=T'T\) in the inner-product computation that makes~\eqref{eq:def}
well-defined).

\item \textbf{Invariance of \(H\)} \(\Rightarrow\) block shape
\(\begin{pmatrix}T'&X\\ 0&Y\end{pmatrix}\).

\item \textbf{Evaluate on the Halmos past} \(\{0\}\oplus\ker T^*\) via \(U^*\)
\(\Rightarrow\)
\[
Xd=T^*T'd,\qquad Yd=DT'd.
\]

\item \textbf{Full commutation of blocks} \(\Rightarrow\)
\[
YT^*=T^*Y,\qquad T'D-XT^*=TX+DY.
\]

\item \textbf{Isometry of blocks} \(\Rightarrow\)
\[
{T'}^*X=0,\qquad X^*X+Y^*Y=I.
\]
\end{enumerate}

The operators \(X,Y\) are uniquely determined by~\eqref{eq:def}; the displayed
equations are simply that construction rewritten in Halmos coordinates.

\bigskip
\begin{remark}[contrast with double commutativity]
If one \emph{also} had \(T^*T'=T'T^*\), then \(T'd\in\ker T^*\) for \(d\in\ker T^*\),
hence \(Xd=T^*T'd=0\) and \(Yd=T'd\). In that case one may take
\[
X=D_{{T'}^*},\qquad Y=-{T'}^*,
\]
i.e.\ \(\widetilde{T}'\) itself is the Halmos matrix of \(T'\). For a merely commuting
pair this shortcut fails, and \(X,Y\) remain the (generally different) operators
cut out by~\eqref{eq:def}.
\end{remark}

\end{document}
